\section{自动反馈与质量保证}
%% ============================================================

让 Claude 自己验证自己的输出是提升质量最有效的方式。

\subsection{给 Claude 一个反馈循环（Tip \#4）}

\begin{importantbox}{2--3x 质量提升}
在 prompt 中包含测试命令、linter 检查或预期输出。Claude 运行测试、看到失败、自动修复，无需你介入。Claude Code 核心开发者 Boris Cherny 表示，仅此一项就能带来 2--3 倍的质量提升。
\end{importantbox}

对于 UI 变更，配置 Playwright MCP server，让 Claude 能打开浏览器、与页面交互、验证 UI 是否正常工作。

\subsection{安装语言级 Code Intelligence 插件（Tip \#5）}

LSP 插件让 Claude 在每次文件编辑后自动获得诊断信息：类型错误、未使用的 import、缺失的返回类型。Claude 在你注意到之前就能看到并修复问题。

\begin{knowledgebox}{这是影响最大的单个插件}
运行 \texttt{/plugin} 进入 Discover 标签浏览完整列表。支持 TypeScript、Python、Rust、Go、C\#、Java、Kotlin、Swift、PHP 等语言。需要系统上安装对应的 Language Server。
\end{knowledgebox}

\subsection{用 gh CLI 和其他命令行工具（Tip \#6）}

\texttt{gh} CLI 处理 PR、Issue、评论，不需要额外的 MCP server。CLI 工具比 MCP server 更节省上下文，因为不需要加载 tool schema。

对于 Claude 不认识的工具，告诉它："Use \texttt{sentry-cli --help} to learn about it, then use it to find the most recent error in production." Claude 会读 help 输出、理解语法、运行命令。

\subsection{对话式 PR Review（Tip \#46）}

不要让 Claude 做"一次性" PR review。在会话中与 Claude 对话式审查：
\begin{itemize}
  \item "Walk me through the riskiest change in this PR."
  \item "What would break if this runs concurrently?"
  \item "Is the error handling consistent with the rest of the codebase?"
\end{itemize}

对话式审查能发现更多问题，因为你可以深入到重要区域。一次性审查往往只标记风格问题，容易遗漏架构问题。

\subsection{本章小结}

质量保证的核心策略：在 prompt 中嵌入反馈循环（测试、lint）、安装 LSP 插件获取实时诊断、用 CLI 工具代替 MCP server 节省上下文、用对话式方法做 PR review。

%% ============================================================
